\scanpage{38}
\margin{\figure{\begin{tikzpicture}[scale=.4]\draw[gray!50] (0,0) .. controls (.5,.1) and (.6,.6) .. (1.5,.55);\draw[gray!30] (0,-.35)--(.12,-.23);\end{tikzpicture}}}
\begin{center}\uline{Preface}\end{center}

The purpose of this book is to study two related problems: when does an orbit space of an algebraic scheme acted on by an algebraic group exist? And to construct moduli schemes for various types of algebraic objects. The second problem appears to be, in essence, a special and highly non-trivial case of the first problem. From an Italian point of view, the crux of both problems is in passing from a birational to a biregular point of view. To construct both orbit spaces and moduli ``generically'' are simple exercises. The problem is whether, within the set of all models of the resulting birational class, there is one model whose geometric points classify the set of orbits in some action, or the set of algebraic objects in some moduli problem. In both cases, it is quite possible that some orbits, or some objects are so exceptional, or, as we shall say, are \uline{not} \uline{stable}, so that they must be left out of the model. The difficulty is to pin down the meaning of stability in a given case. One of the most intriguing unsolved problems, in this regard, is that of the moduli of non-singular polarized surfaces. Which such surfaces are not stable, in the sense that there is no moduli scheme for them and their deformations? This property is very delicate.

One of my principles has been not to worry too much about the difference between characteristic $0$, and finite characteristics. A large part of this book is, therefore, devoted to a theory developed only in characteristic $0$. I am convinced, however, that it is almost entirely valid in all characteristics. What is necessary is to find

\mspage{37}
